\section{Adversarial inputs for heuristic hashes}
\label{app:adversarial}

This appendix records the inputs behind the heuristic rows of
\Cref{tab:injective:adversarial}: for each hash, the message pair we used, the
measured collision rate, and the number of random secrets it was measured over.
All code is in \texttt{tools/bench/adversarial/} of the accompanying repository
(self-contained C++17; every competitor hash is re-implemented from its
published algorithm and checked against published test vectors, e.g.\
\texttt{XXH3} against the official sanity vectors for lengths
$12, 24, 48, 80, 195$ with seed $0$ and seed \texttt{PRIME64}).  Measurements
were taken on an Apple M2 Pro with Apple clang~17 at \texttt{-O3}.

\paragraph{Model and methodology.}
As in \Cref{sec:injective:experiments}, the secret (seed or secret array) is
drawn uniformly at random per trial and hidden; the attacker chooses a
\emph{fixed} pair of messages; the reported quantity is the fraction of secrets
for which the pair collides.  A random function gives $2^{-64}$.  Attacks that
need a hash's \emph{public} default constants (the shipped \texttt{wyhash} and
\texttt{rapidhash} secret arrays, \texttt{XXH3}'s \texttt{kSecret} with seed
$0$) do not count in this model and are listed separately at the end.  A
message is a byte string whose length is a multiple of $16$, read as
little-endian 64-bit words $w_0,w_1,\ldots$; $M$ denotes the all-ones word, and
``complementing'' a word means XORing it with $M$.  Heuristic hashes were run
at $2^{29}$ trials per cell in the site scan and $2^{31}$ trials per cell in the
length scan; the proven baselines at $2^{27}$ and $2^{29}$ respectively.

\paragraph{The multiply--fold differential.}
All four heuristic hashes share a $64\times64\to128$-bit multiply whose two
halves are folded together, $\mathrm{lo}_{64}(uQ)\oplus\mathrm{hi}_{64}(uQ)$
(\texttt{wyhash}, \texttt{rapidhash}, \texttt{XXH3}) or
$\mathrm{lo}_{64}(uQ)+\mathrm{hi}_{64}(uQ)$ (\texttt{MUM}).  The operands
$u,Q$ are a message word combined with a secret or state word, so an attacker
controls only the difference $(\delta,\varepsilon)$ that the two messages impose
on $(u,Q)$, while $(u,Q)$ themselves are uniform.  The relevant quantity is
therefore $P(\delta,\varepsilon)=\Pr_{u,Q}[\mathrm{fold}(u,Q)=
\mathrm{fold}(u\oplus\delta,Q\oplus\varepsilon)]$, which
\texttt{fold\_search.cpp} evaluates exhaustively for $8$- and $16$-bit operands
and by $2^{34}$ samples at $24$--$64$ bits.  For the XOR fold the best
differences are $(M,M)$ and $(M,0)$: at $64$ bits they collide with probability
$2^{-26.7}$ and $2^{-26.8}$, and the exponent scales as $P\approx2^{-0.42n}$ in
the operand width $n$ rather than the ideal $2^{-n}$.  For the additive fold of
\texttt{MUM} the difference $(M,M)$ collides with probability $0.667$ and the
additive difference $(0x55\cdots5,0)$ with probability $0.185$, both
\emph{independent of the operand width}.  Complementing an operand is the single
most powerful difference we found.

\paragraph{\texttt{wyhash} and \texttt{rapidhash} (random secret).}
Versions: \texttt{wyhash} final v4.3 (with \texttt{WYHASH\_\allowbreak CONDOM=1}) and
\texttt{rapidhash} v1.0, both with a uniformly random secret array in place of
the shipped one.  The first step of both hashes is
$\mathrm{mix}(w_0\oplus\mathrm{secret}, w_1\oplus\mathrm{state})$, so
complementing $w_0$ and $w_1$ imposes the $(M,M)$ difference on the first fold.
Measured over $2^{31}$ random secrets per length, the pair collides with
probability $2^{-26.6}$, $2^{-26.7}$, $2^{-26.2}$, $2^{-26.6}$, $2^{-26.4}$
(\texttt{wyhash}) and $2^{-26.4}$, $2^{-26.5}$, $2^{-26.5}$, $2^{-26.2}$,
$2^{-26.8}$ (\texttt{rapidhash}) at $32$, $48$, $64$, $96$ and $160$ bytes: flat
in the message length and unaffected by the random secret, as the
fold analysis predicts.  Complementing a single interior word instead
propagates through two folds and drops to the noise floor of these trial counts.

\paragraph{\texttt{XXH3} (random seed).}
\texttt{XXH3\_64bits\_withSeed} on the $9$--$240$-byte paths, seed drawn at
random per trial.  Complementing $w_0$ and $w_1$ (the operands of the first
\texttt{mix16B}) collides with probability $2^{-26.3}$, $2^{-25.8}$,
$2^{-26.1}$, $2^{-26.2}$, $2^{-26.2}$ at $32$, $48$, $64$, $96$, $160$ bytes
($2^{31}$ seeds per length).  Single-word complements of $w_1$, $w_2$, $w_3$ in
a $32$-byte message give $2^{-26.4}$, $2^{-27.0}$, $2^{-28.0}$ ($2^{29}$
seeds; these three rest on only $6$, $4$ and $2$ observed collisions, so they
locate the rate only to within a factor of about two).  A re-test against the
newest upstream release (v0.8.3, byte-identical to the packaged header) and the
development branch, with $2^{30}$ seeds per cell, reproduced these rates and
sharpened the picture: for the $64$-bit variant the effect lives entirely in
the multiply--fold, so complementing $w_0$ alone, or both words with an
unrelated $w_1$, collides at comparable rates, and the rate then depends on the
rest of the message; one $128$-byte pair reached $2^{-23.3}$ ($97$ and $109$
collisions in two independent runs of $2^{30}$ seeds), which is the entry in
\Cref{tab:injective:adversarial}.  A uniformly random $192$-byte secret in
place of the seed changes nothing.  We reported the differential upstream as
xxHash issue~\#1127.\footnote{\url{https://github.com/Cyan4973/xxHash/issues/1127}}

\paragraph{\texttt{XXH3-128}.}
The $128$-bit variant does not escape the differential.  On its $17$--$128$-byte
path each of the two accumulators adds the multiply--fold of one $16$-byte
chunk and XORs in the \emph{raw word sum} of the other chunk.  Choosing
$w_1=\lnot w_0$ makes that sum invariant under complementing both words, so
one accumulator is unchanged deterministically and the other collides exactly
when the fold does.  Measured over $2^{29}$ random seeds, the full $128$-bit
output collides with probability $2^{-27.0}$, $2^{-27.0}$, $2^{-26.2}$,
$2^{-26.0}$, $2^{-26.4}$ and $2^{-26.0}$ at $32$, $48$, $64$, $100$, $128$ and
$160$ bytes, every collision being a joint collision of both halves;
complementing $w_0$ alone, or both words with an unrelated $w_1$, gave none.
Doubling the output width therefore buys nothing against this input pair, and
\texttt{XXH3-128} scores $28.5$ bits in \Cref{tab:injective:adversarial}, barely above the $27.3$ of \texttt{XXH3}.

\paragraph{\texttt{komihash} (random seed).}
\texttt{komihash} v5.34 (checked against its $63$ published test vectors)
resisted every differential we tried: the $(M,M)$ and $(M,0)$ complements on
the fold operands of its minimal-diffusion path for $8$--$15$-byte messages,
single- and double-word complements in every $16$-byte block position,
length-boundary and zero-padding pairs, and the structural pairs used against
\texttt{MUM}.  Every surface gave $0$ collisions in $2^{32}$ random seeds,
consistent with the $2^{-64}$ baseline.  We therefore list it as ``$\ge32$'' in
\Cref{tab:injective:adversarial}: the search shows only that no pair we found
exceeds $2^{-32}$, and is not a proof.

\paragraph{\texttt{MUM} v3.}
Upstream \texttt{mum.h} (re-implemented in \texttt{hashes.h}) with its default macros, in both the unroll-8 (x86-64) and
unroll-16 (aarch64) configurations.  The additive fold satisfies
$\_\mathrm{mum}(v,p)+\_\mathrm{mum}(\lnot v,p)\in\{2^{64}-1,\,2^{64}-2\}$ (mod~$2^{64}$)
for every $v$ (the second value occurs exactly when the low halves of $vp$ and
$(\lnot v)p$ carry, i.e.\ for a $p/2^{64}$ fraction of the $v$---between $0.15$
and $0.87$ for \texttt{MUM}'s primes; \texttt{selftest.cpp} checks the analogous
identity for complementing both operands), so complementing an input
word toggles that word's contribution to the XOR accumulator by
approximately~$M$.  Complementing any single interior word of a $32$-byte
message therefore collides with probability $2^{-1}$ (the finalizer maps the
flipped and unflipped states together about half the time), and complementing
\emph{two adjacent} interior words $w_1,w_2$ cancels the two toggles and
collides for \emph{every} seed: $2^{-0.0}$ over $2^{29}$ seeds, i.e.\ a
key-free collision.  The shortest key-free pair we know is a single $8$-byte
word: $w=\texttt{0xb0899b7198a95479}$ and $w'=\texttt{0x08eb9f259b16692d}$
satisfy $\_\mathrm{mum}(w,p_0)=\_\mathrm{mum}(w',p_0)$ (found by a Brent
$\rho$ search on $v\mapsto\_\mathrm{mum}(v,p_0)$ and checked over $2^{20}$
seeds), which is what places \texttt{MUM} at $0$ bits in
\Cref{tab:injective:adversarial}.  Two further seed-independent constructions were found: a
pair of values $w_1$ that collide under $\_\mathrm{mum}(\cdot,p_1)$ (a
$\rho$-collision located after $2^{37}$ evaluations), and, at $160$ bytes, a
rotation clique $w_0=p_0\oplus2^k$, $w_1=\mathrm{rotr}(Q,k)\oplus p_1$ of
$K=64$ messages that are pairwise colliding for all $1020$ seeds tried.

\paragraph{The paper's own fold variant.}
For completeness, the recurrence of \Cref{sec:injective} with the field product
replaced by the MUM fold (the throughput data point of
\Cref{sec:injective:experiments}) is subject to the same differential:
complementing the last $b$-word imposes $(M,0)$ on the final fold and collides
with probability $2^{-27.5}$, $2^{-26.7}$, $2^{-27.1}$, $2^{-27.3}$,
$2^{-27.1}$ at $32$--$160$ bytes ($2^{31}$ keys).  This is the variant the text
explicitly excludes from the Schwartz--Zippel guarantee.

\paragraph{Shipped defaults (excluded from the model).}
With the \emph{published} constants, forced collisions exist: placing
$\mathrm{secret}[1]$ at byte offset $\mathrm{len}-16$ annihilates the final
multiplication of \texttt{wyhash} and \texttt{rapidhash} (rate $1.000$ over
$2^{20}$ seeds, at $32$ and $160$ bytes); $w_0=\mathrm{secret}[1]$
(\texttt{wyhash}) or $w_0=\mathrm{secret}[2]$ (\texttt{rapidhash}) zeroes the
block state so that messages differing only in $w_1$ collide; and for
unseeded \texttt{XXH3}, $w_0=\texttt{kSecret}[0..8)$ zeroes
\texttt{mix16B} of the first chunk.  Each of these yields a set of $K=256$
messages that collide pairwise for all $1020$ seeds tried.  A random secret
removes all of them (rate $0$ in the same experiments), which is why they are
excluded from \Cref{tab:injective:adversarial}; \texttt{MUM}'s constructions
have no secret to randomize and remain.

\paragraph{Proven baselines and controls.}
On every differential site, length, annihilation attempt and blocking set
above, the field version of our recurrence over $\F_{2^{64}}$, its Mersenne
variant, and Dietzfelbinger's vector multiply-shift produced $0$ collisions
($2^{24}$--$2^{29}$ trials per cell, and $1020$ seeds for the blocking sets), consistent with their $d/|\F|$ and
$2^{-64}$ bounds.  On two \emph{uniformly random} messages every hash,
heuristic or proven, produced $0$ collisions at $32$ and $160$ bytes
($2^{30}$ trials for the heuristic hashes, $2^{28}$ for the proven ones): the
elevated rates above are entirely input-specific.

\paragraph{A seeded-differential test for SMHasher3.}
SMHasher3 runs every hash under fixed seeds, so the input-specific rates above
are invisible to it: they are probabilities over the \emph{seed} for a fixed
pair of inputs.  We added a test (\texttt{SeedDifferential}, in the paper's
fork of SMHasher3) that draws $2^{24}$ uniformly random seeds ($2^{30}$ in an
extended tier) and counts full-output collisions of fixed structured pairs
derived from a base message whose second word is the complement of its first:
complementing the first two words, single words, or adjacent interior pairs,
at lengths from $16$ bytes to $1$\,KB.  A row fails when its count exceeds
$\max(3,\,E+6\sqrt{E})$ for the expected count $E=N2^{-w}$.  In the extended
tier \texttt{wyhash}, \texttt{rapidhash}, \texttt{XXH3} and \texttt{XXH3-128}
fail at about $2^{-26}$ on the complemented leading pair and every
\texttt{MUM} variant fails on the adjacent pairs at rate $1$, while
\texttt{komihash}, Polymur, SipHash, \texttt{t1ha2}, FarmHash, CityHash, the
Mersenne polynomial hash and ChainHash show no collision in $2^{30}$ seeds.

\paragraph{Caveat.}
The rates reported here are for the worst inputs \emph{we found}, by a
differential search restricted to word complements and additive differences on
the fold operands, together with one-line annihilation constructions.  They are
lower bounds on the true worst-case collision probability of each hash; inputs
with a higher rate may exist.  Each $2^{-26}$--$2^{-27}$ entry rests on between
$11$ and $36$ observed collisions in $2^{31}$ trials, so its exponent carries a
statistical uncertainty of about $\pm0.3$ ($\pm0.4$ at the smallest counts); the differences between lengths and
between the three XOR-fold hashes are within that noise.  What the experiments
establish is a lower bound of order $2^{-27}$ (\texttt{wyhash},
\texttt{rapidhash}, \texttt{XXH3}) and $1$ (\texttt{MUM}) on the worst case,
against a proven upper bound of $n/2^{64}$ for the three-key field version of
our hash on $n$ word pairs ($(3n-1)/2^{64}$ with a single key word).

\paragraph{Truncated Mersenne outputs.}
A polynomial hash over $\F_p$ with $p=2^{89}-1$ must not be shortened by
keeping the low $64$ bits of the residue.  Two messages that differ only in
the final additive word, by $2^{64}-1$, have residues that differ by exactly
that constant; whenever the sum wraps past $p$, the wrap subtracts
$p\equiv-1\pmod{2^{64}}$ and the low $64$ bits coincide.  This happens
whenever the residue lies in the top $2^{64}$ values of $[0,p)$, so, over a
random key, the pair collides with probability about $2^{64}/p=2^{-25}$,
whatever the remaining words of the message (we measured $2$ collisions in
$2^{26}$ random keys and $8$ in $2^{28}$).  The same mechanism gives
$2^{32}/(2^{61}-1)=2^{-29}$ for the low $32$ bits of a residue modulo
$2^{61}-1$ (measured $2$ in $2^{30}$ and $8$ in $2^{32}$).  The Mersenne rows of
\Cref{tab:injective:adversarial} therefore output the full residue; a
shorter output needs a universal post-hash, for instance one multiply-shift
step or a degree-$k$ chain as in \Cref{sec:main-theorem}.
